% Cette fiche d'activité à été développée par
% + Erwan Tanguy-Legac <erwan.tanguy-legac@ens-rennes.fr>
% + Héloïse Troublé <heloise.trouble@ens-rennes.fr>

% Le template a été largement inspiré de celui de
% + Santiago Bautista <Santiago.Bautista@ens-rennes.fr>
% + Benjamin Bordais <Benjamin.Bordais@ens-rennes.fr>

% Elle est distribuée sous la licence Creative-Commons by-nc-sa 4.0
% (Attribution, Non-Commercial, Share Alike)
% qui peut être trouvée dans
% [le site de Creative Commons](https://creativecommons.org/licenses/by-nc-sa/4.0/)

\documentclass[12pt, a4paper]{article}
\usepackage[french]{babel}
\usepackage[utf8]{inputenc}
\usepackage[top=0.9in, left=0.6in, right=0.7in, bottom=1in]{geometry}
\usepackage{amssymb}
\usepackage{color}
\usepackage{fancyhdr}
\usepackage{longtable}
\usepackage{array}
\usepackage{titling}

\pagestyle{fancy}

\renewcommand{\headrulewidth}{1pt}
\chead{}
\lhead{Marmottes - Plan de l'activité}
\rhead{Héloïse Troublé et Erwan Tanguy-Legac}

\renewcommand{\footrulewidth}{1pt}
\cfoot{\thepage} 
\lfoot{2024}
\rfoot{ENS Rennes}

\fancypagestyle{plain}{
	\fancyhf{} 
	\cfoot{\thepage}
	\renewcommand{\headrulewidth}{0pt}
	\renewcommand{\footrulewidth}{1pt}}

\begin{document}
	
	\title{Marmottes}
	\author{Héloïse Troublé et Erwan Tanguy-Legac}
	\date{29 mars 2024}
	\maketitle
	
	\paragraph*{Niveau :} $6^{\grave{e}me}$
	\paragraph*{Durée :} 55 minutes
	\paragraph*{Prérequis :} Aucun
	\paragraph*{Description du jeu :}
  L'activité des marmottes consiste à introduire l'algorithme de compression de Huffman, décrit ci-dessous, et à présenter l'encodage
  d'un texte selon un arbre de Huffman.
	
	\paragraph*{Objectif pédagogique :}
  Découverte de l'algorithme de Huffman, introduction aux arbres binaires, introduction à la compression sans perte.

  \paragraph*{Matériel nécéssaire :}
  Pour chaque groupe de 4 élèves, prévoir minimum 4 fourches, 5 marmottes (4, 3, 2, 2, 1), et une feuille pour la partie décodage.
  Il faudrait prévoir des instances supplémentaires pour les groupes les plus rapides (une fourche et une marmotte 1 en plus d'une deuxième
  instance de décodage suffisent).
	
	\begin{longtable}{c|m{3.9cm}|m{8.4cm}|m{2.4cm}}		\textbf{Durée} & \textbf{Phase} & \textbf{Activités et consignes} & \textbf{Organisation et matériel} \\ \hline\hline
		
		10' &
		
		Introduction
		
		&
		
		On se présente, on explique les règles clairement (surtout pour le comptage des points), on s'assure le plus possible que tout
    le monde ait compris.
		
		&
		
		A l'oral au tableau (avec vidéoprojecteur ou aimants selon la salle)
		
		\\ \hline
		
	  10' &
		
    Recherche
		
		&
		
	  On laisse les élèves chercher un arbre optimal, on passe dans les rangs pour s'assurer que tout le monde ait bien compris les règles.
    Pour les plus rapides, on leur donne d'autres instances.
		
		& Par groupes de 4 (en îlots)
		
		\\ \hline
		
		5' + 5' &
		
		Mise en commun et explication de la stratégie optimale
		
		&
		
	  On demande aux élèves comment ils ont fait, puis on leur présente l'algorithme de Huffman. L'idéal est d'avoir réussi
    à trouver un groupe qui a bien avancé et leur expliquer rapidement l'algorithme en leur demandant s'ils acceptent de
    passer au tableau pour le présenter au reste de la classe.

		&
		
		Au tableau (avec vidéoprojecteur ou aimants selon la salle)
		
		\\ \hline
		
		5'
		
		&
		
	  Présentation décodage
		
		&
		
		On présente la méthode de décodage d'un texte qui a été encodé avec un arbre de Huffman.
		
		& Idem
		
		\\ \hline
		
		15' &
		
	  Décodage et encodage
		
		&
		
		On laisse les élèves décoder un texte qu'on leur a distribué (avec l'arbre associé). Ensuite si la durée le permet,
    on leur propose de se séparer en groupes binômes et d'encoder un mot avant de le transmettre à l'autre binôme du groupe
    de 4. Les élèves sont très lents sur l'encodage.
		
		&
		
		Par groupes de 4
		
		\\ \hline
		
		5'
		
		&
		
		Institutionnalisation
		
		&
		
		On reprend l'attention des élèves, et on présente l'institutionnalisation (détail ci-dessous). Les grands axes sont la
    présentation du stockage des données dans un ordinateur, une rapide explication de l'algorithme de construction, puis
    les cas d'utilisation de la compression.
		
		&
		
		A l'oral
		
	\end{longtable}
	
	
\subsection*{Présentation :}
C’est l’hiver, les marmottes s’apprêtent à hiberner. Elles veulent pour cela construire un terrier avec plein de galeries. Afin de
s’assurer que le terrier ne s’effondre pas, on ne peut que construire des galeries en forme de V inversé (il est très important ici
d'expliciter la structure de l'arbre, avec un parent et deux fils). Mais on a un problème : de temps en temps, elles se réveillent
et remontent à la surface, en faisant plein de bruit, ce qui risquerait de réveiller les autres. À cause de ça, aucune marmotte ne
veut dormir au milieu d’une galerie : elles dorment toutes au fond (sinon elles se feraient marcher dessus). On veut construire une
galerie pour que les marmottes fassent le moins de bruit possible. Pour déterminer le bruit que fait une marmotte, on compte le nombre
de galeries qu'elle doit traverser pour remonter à la surface (ce qui correspond à sa profondeur dans l'arbre). Le chiffre sur chaque
marmotte indique le nombre de fois qu’une marmotte va se réveiller au cours de l’hiver. (petit exemple)
On va maintenant vous distribuer de quoi faire des galeries, le but est d’en construire une telle que le moins de marmottes possibles
soient réveillées au cours de l’hiver.

\subsection*{Présentation de l'algorithme :}
Il existe une méthode pour construire à tous les coups la meilleure galerie. Pour cela, il faut commencer par ordonner les marmottes
en fonction du nombre de fois qu’elles se réveillent. On prend ensuite les deux premières, et on les met toutes les deux à l’extrémité
d’une galerie. On compte le nombre de réveils qui auront lieu dans cette galerie, et on place cette galerie avec les marmottes, en
gardant l’ordre qu’on avait. On répète cela tant qu’on a deux galeries. On obtient ainsi à la fin une très grande galerie, qui
est la meilleure possible.

\subsection*{Présentation décodage :}
Ce que l’on vient de faire sert à raccourcir la taille des textes dans l’ordinateur. Les lettres sont stockées comme une suite de 8 1 ou
de 0 (comme toute information). Ce n’est pas toujours très efficace donc on aimerait réduire la place que ça prend. Si on remplace les
marmottes par les lettres, et les couloirs par un choix entre 0 et 1, on peut coder et décoder un texte (c'est comme un code secret).
On va vous faire passer une galerie construite avec la méthode que vous venez de voir, mais on a remplacé les marmottes par des lettres.
On va en plus vous donner une suite de 0 et de 1, et vous devrez trouver la phrase qui a été codée. Pour décoder une lettre, vous partez
de l’entrée de la galerie et si le prochain chiffre est un 0, vous allez à gauche, sinon vous allez à droite. Quand vous arrivez sur une
lettre, vous écrivez la lettre et vous repartez de l’entrée de la galerie pour décoder la lettre suivante. (Petit exemple avec un
arbre de hauteur 2 au tableau, le mot OUI fonctionne bien)

\subsection*{Institutionnalisation :}
C’est de l’informatique, parce que grâce au codage des lettres créé par le terrier, on a réussi à coder un texte en prenant très peu de place !
En informatique, c’est important dans la gestion de la mémoire de l’ordinateur : on veut pouvoir stocker un maximum d’informations, de données
dans un minimum de place. Pour ça, on réfléchit à une façon intelligente d’encoder les informations qui sont obligatoirement stockées dans
l’ordinateur comme une suite de 0 et de 1.

Ce qu’on a appliqué ici, c’est l’encodage d’Huffman. On regarde la fréquence d’apparition des caractères dans un texte, et on applique
l’algorithme d’Huffman sur ces fréquences. (ex : si une lettre est présente 3 fois dans le texte, on ajoute une marmotte 3). Ensuite on
construit le terrier qui minimisera le nombre de fois où des marmottes se réveillent, en fusionnant les plus petits terriers. Ça construit
un codage de chaque caractère qui fait en sorte qu’au final, le stockage du texte prenne le moins de place possible.

Faire en sorte que les données prennent moins de place dans l’ordinateur, c’est super important. C’est ça qui permet de voir des vidéos
en direct et de stocker plein de photos sur une clé USB. Vous savez qu’on peut créer toutes les couleurs avec les 3 couleurs primaires
(leur demander lesquelles c’est, leur dire qu’en informatique c’est rouge/vert/bleu). On peut voir une image comme une grille, où chaque
case contient 3 lettres pour indiquer la quantité de chaque couleur, ça fait beaucoup de lettres à stocker (sachant que les grilles sont
de taille 1000*1000). Compresser permet de stocker environ 250 fois plus d’images, ce qui est beaucoup.

\subsection*{Instances :}
On commence avec une instance assez simple, composée des jetons 4 3 2 2 1. Cette instance a pour optimal une valeur de 27, avec le 4 à une
profondeur 2 (le 3 peut soit être avec le 4 à profondeur 2 ou à profondeur 3). Rajouter une marmotte 1 à cette instance change la
valeur de l'optimal à 32.

Deux instances pour le décodage sont fournies dans le répertoire correspondant (une courte et une plus longue).

\end{document}
